\documentclass[12pt]{article}

% Idioma y codificación
\usepackage[spanish, es-tabla, es-notilde]{babel}       %es-tabla para que se titule "Tabla"
\usepackage[utf8]{inputenc}

% Márgenes
\usepackage[a4paper,top=3cm,bottom=2.5cm,left=3cm,right=3cm]{geometry}

% Comentarios de bloque
\usepackage{verbatim}

% Paquetes de links
\usepackage[hidelinks]{hyperref}    % Permite enlaces
\usepackage{url}                    % redirecciona a la web

% Más opciones para enumeraciones
\usepackage{enumitem}

% Personalizar la portada
\usepackage{titling}

% Paquetes de tablas
\usepackage{multirow}

% Para añadir el símbolo de euro
\usepackage{eurosym}

% Paquete para subrayar con diferentes estilos
\usepackage[normalem]{ulem}


%------------------------------------------------------------------------

%Paquetes de figuras
\usepackage{caption}
\usepackage{subcaption} % Figuras al lado de otras
\usepackage{float}      % Poner figuras en el sitio indicado H.


% Paquetes de imágenes
\usepackage{graphicx}       % Paquete para añadir imágenes
\usepackage{transparent}    % Para manejar la opacidad de las figuras

% Paquete para usar colores
\usepackage[dvipsnames, table, xcdraw]{xcolor}
\usepackage{pagecolor}      % Para cambiar el color de la página

% Habilita tamaños de fuente mayores
\usepackage{fix-cm}

% Para los gráficos
\usepackage{tikz}
\usepackage{forest}

% Para poder situar los nodos en los grafos
\usetikzlibrary{positioning}

% Para evitar problemas entre babel y tikz
\usetikzlibrary{babel}


%------------------------------------------------------------------------

% Paquetes de matemáticas
\usepackage{mathtools, amsfonts, amssymb, mathrsfs}
\usepackage[makeroom]{cancel}     % Simplificar tachando
\usepackage{polynom}    % Divisiones y Ruffini
\usepackage{units} % Para poner fracciones diagonales con \nicefrac

\usepackage{pgfplots}   %Representar funciones
\pgfplotsset{compat=1.18}  % Versión 1.18

\usepackage{tikz-cd}    % Para usar diagramas de composiciones
\usetikzlibrary{calc}   % Para usar cálculo de coordenadas en tikz

%Definición de teoremas, etc.
\usepackage{amsthm}
%\swapnumbers   % Intercambia la posición del texto y de la numeración

\theoremstyle{plain}

\makeatletter
\@ifclassloaded{article}{
  % --- Configuración para ARTICLE ---
  \newtheorem{teo}{Teorema}[section]
  \theoremstyle{definition}
  \newtheorem{definicion}{Definición}[section]
  \newtheorem{ejercicio}{Ejercicio}[section]

  \renewcommand{\theteo}{%
    \ifnum\value{section}=0 \arabic{teo}\else \thesection.\arabic{teo}\fi
  }
  \renewcommand{\thedefinicion}{%
    \ifnum\value{section}=0 \arabic{definicion}\else \thesection.\arabic{definicion}\fi
  }
  \renewcommand{\theejercicio}{%
    \ifnum\value{section}=0 \arabic{ejercicio}\else \thesection.\arabic{ejercicio}\fi
  }
}{
  % --- Configuración para BOOK / REPORT ---
  \newtheorem{teo}{Teorema}[chapter]
  \theoremstyle{definition}
  \newtheorem{definicion}{Definición}[chapter]
  \newtheorem{ejercicio}{Ejercicio}[section]

  % Si no hay sección: Capítulo.Teorema (ej. 2.1)
  % Si hay sección:    Capítulo.Sección.Teorema (ej. 2.3.1)
  \renewcommand{\theteo}{%
    \ifnum\value{section}=0
      \thechapter.\arabic{teo}%
    \else
      \thesection.\arabic{teo}% (\thesection ya contiene el número de capítulo)
    \fi
  }
  \renewcommand{\thedefinicion}{%
    \ifnum\value{section}=0
      \thechapter.\arabic{definicion}%
    \else
      \thesection.\arabic{definicion}%
    \fi
  }
  \renewcommand{\theejercicio}{%
    \ifnum\value{section}=0
      \arabic{ejercicio}%
    \else
      \thesection.\arabic{ejercicio}% Formato por defecto: ej. 2.2.1
    \fi
  }
}
\makeatother

% Teoremas dependientes o con contador compartido
\newtheorem{coro}{Corolario}[teo]
\newtheorem{prop}[teo]{Proposición}
\newtheorem{lema}[teo]{Lema}
\newtheorem*{lema*}{Lema}

\theoremstyle{remark}
\newtheorem*{observacion}{Observación}

\theoremstyle{definition}
\newtheorem*{notacion}{Notación}
\newtheorem*{ejemplo}{Ejemplo}
\newtheorem*{ejercicio*}{Ejercicio}


% \renewcommand\qedsymbol{$\blacksquare$}         % Cambiar símbolo QED
%------------------------------------------------------------------------

% Paquetes para encabezados
\usepackage{fancyhdr}
\pagestyle{fancy}
\fancyhf{}

\newcommand{\helv}{ % Modificación tamaño de letra
\fontfamily{}\fontsize{12}{12}\selectfont}
\setlength{\headheight}{15pt} % Amplía el tamaño del índice


%\usepackage{lastpage}   % Referenciar última pag   \pageref{LastPage}
%\fancyfoot[C]{%
%  \begin{minipage}{\textwidth}
%    \centering
%    ~\\
%    \thepage\\
%    \href{https://losdeldgiim.github.io/}{\texttt{\footnotesize losdeldgiim.github.io}}
%  \end{minipage}
%}
\fancyfoot[C]{\thepage}
\fancyfoot[R]{\href{https://losdeldgiim.github.io/}{\texttt{\footnotesize losdeldgiim.github.io}}}

%------------------------------------------------------------------------

% Conseguir que no ponga "Capítulo 1". Sino solo "1."
\makeatletter
\@ifclassloaded{book}{
  \renewcommand{\chaptermark}[1]{\markboth{\thechapter.\ #1}{}} % En el encabezado
    
  \renewcommand{\@makechapterhead}[1]{%
  \vspace*{50\p@}%
  {\parindent \z@ \raggedright \normalfont
    \ifnum \c@secnumdepth >\m@ne
      \huge\bfseries \thechapter.\hspace{1em}\ignorespaces
    \fi
    \interlinepenalty\@M
    \Huge \bfseries #1\par\nobreak
    \vskip 40\p@
  }}
}
\makeatother

%------------------------------------------------------------------------
% Paquetes de cógido
\usepackage{minted}
\renewcommand\listingscaption{Código fuente}

\usepackage{fancyvrb}
% Personaliza el tamaño de los números de línea
\renewcommand{\theFancyVerbLine}{\small\arabic{FancyVerbLine}}

% Estilo para C++
\newminted{cpp}{
    frame=lines,
    framesep=2mm,
    baselinestretch=1.2,
    linenos,
    escapeinside=||
}

% para minted
\definecolor{LightGray}{rgb}{0.95,0.95,0.92}
\setminted{
    linenos=true,
    stepnumber=5,
    numberfirstline=true,
    autogobble,
    breaklines=true,
    breakautoindent=true,
    breaksymbolleft=,
    breaksymbolright=,
    breaksymbolindentleft=0pt,
    breaksymbolindentright=0pt,
    breaksymbolsepleft=0pt,
    breaksymbolsepright=0pt,
    fontsize=\footnotesize,
    bgcolor=LightGray,
    numbersep=10pt
}


\usepackage{listings} % Para incluir código desde un archivo

\renewcommand\lstlistingname{Código Fuente}
\renewcommand\lstlistlistingname{Índice de Códigos Fuente}

% Definir colores
\definecolor{vscodepurple}{rgb}{0.5,0,0.5}
\definecolor{vscodeblue}{rgb}{0,0,0.8}
\definecolor{vscodegreen}{rgb}{0,0.5,0}
\definecolor{vscodegray}{rgb}{0.5,0.5,0.5}
\definecolor{vscodebackground}{rgb}{0.97,0.97,0.97}
\definecolor{vscodelightgray}{rgb}{0.9,0.9,0.9}

% Configuración para el estilo de C similar a VSCode
\lstdefinestyle{vscode_C}{
  backgroundcolor=\color{vscodebackground},
  commentstyle=\color{vscodegreen},
  keywordstyle=\color{vscodeblue},
  numberstyle=\tiny\color{vscodegray},
  stringstyle=\color{vscodepurple},
  basicstyle=\scriptsize\ttfamily,
  breakatwhitespace=false,
  breaklines=true,
  captionpos=b,
  keepspaces=true,
  numbers=left,
  numbersep=5pt,
  showspaces=false,
  showstringspaces=false,
  showtabs=false,
  tabsize=2,
  frame=tb,
  framerule=0pt,
  aboveskip=10pt,
  belowskip=10pt,
  xleftmargin=10pt,
  xrightmargin=10pt,
  framexleftmargin=10pt,
  framexrightmargin=10pt,
  framesep=0pt,
  rulecolor=\color{vscodelightgray},
  backgroundcolor=\color{vscodebackground},
}

%------------------------------------------------------------------------
% La etiqueda IA sirve para indicar a partir de donde le pasamos el preambulo a la IA para el escaneo
% IA

% Comandos definidos
\newcommand{\bb}[1]{\mathbb{#1}}
\newcommand{\cc}[1]{\mathcal{#1}}

% I prefer the slanted \leq
\let\oldleq\leq % save them in case they're every wanted
\let\oldgeq\geq
\renewcommand{\leq}{\leqslant}
\renewcommand{\geq}{\geqslant}

% Si y solo si
\newcommand{\sii}{\iff}

% MCD y MCM
\DeclareMathOperator{\mcd}{mcd}
\DeclareMathOperator{\mcm}{mcm}

% Signo
\DeclareMathOperator{\sgn}{sgn}

% Imagen
\DeclareMathOperator{\Ima}{Im}

% Aplicación identidad
\newcommand{\id}{\mathrm{Id}}

% Letras griegas
\newcommand{\eps}{\epsilon}
\newcommand{\veps}{\varepsilon}
\newcommand{\lm}{\lambda}

\newcommand{\ol}{\overline}
\newcommand{\ul}{\underline}
\newcommand{\wt}{\widetilde}
\newcommand{\wh}{\widehat}

\let\oldvec\vec
\renewcommand{\vec}{\overrightarrow}

% Derivadas parciales
\newcommand{\del}[2]{\frac{\partial #1}{\partial #2}}
\newcommand{\Del}[3]{\frac{\partial^{#1} #2}{\partial #3^{#1}}}
\newcommand{\deld}[2]{\dfrac{\partial #1}{\partial #2}}
\newcommand{\Deld}[3]{\dfrac{\partial^{#1} #2}{\partial #3^{#1}}}


\newcommand{\AstIg}{\stackrel{(\ast)}{=}}
\newcommand{\Hop}{\stackrel{L'H\hat{o}pital}{=}}

\newcommand{\red}[1]{{\color{red}#1}} % Para integrales, destacar los cambios.

% Método de integración
\newcommand{\MetInt}[2]{
    \left[\begin{array}{c}
        #1 \\ #2
    \end{array}\right]
}

% Declarar aplicaciones
% 1. Nombre aplicación
% 2. Dominio
% 3. Codominio
% 4. Variable
% 5. Imagen de la variable
\newcommand{\Func}[5]{
    \begin{equation*}
        \begin{array}{rrll}
            \displaystyle #1:& \displaystyle  #2 & \longrightarrow & \displaystyle  #3\\
               & \displaystyle  #4 & \longmapsto & \displaystyle  #5
        \end{array}
    \end{equation*}
}

%------------------------------------------------------------------------


\newcommand{\myhref}[2]{\href{#1}{\color{blue}\uline{#2}}}

% // TODO: Publicar
\begin{document}

    % 1. Foto de fondo
    % 2. Título
    % 3. Encabezado Izquierdo
    % 4. Color de fondo
    % 5. Coord x del titulo
    % 6. Coord y del titulo
    % 7. Fecha

    
    % 1. Foto de fondo
% 2. Título
% 3. Encabezado Izquierdo
% 4. Color de fondo
% 5. Coord x del titulo
% 6. Coord y del titulo
% 7. Fecha
% 8. Autor

\newcommand{\portada}[8]{

    \portadaBase{#1}{#2}{#3}{#4}{#5}{#6}{#7}{#8}
    \portadaBook{#1}{#2}{#3}{#4}{#5}{#6}{#7}{#8}
}

\newcommand{\portadaExamen}[8]{

    \portadaBase{#1}{#2}{#3}{#4}{#5}{#6}{#7}{#8}
    \portadaArticle{#1}{#2}{#3}{#4}{#5}{#6}{#7}{#8}
}




\newcommand{\portadaBase}[8]{

    % Tiene la portada principal y la licencia Creative Commons
    
    % 1. Foto de fondo
    % 2. Título
    % 3. Encabezado Izquierdo
    % 4. Color de fondo
    % 5. Coord x del titulo
    % 6. Coord y del titulo
    % 7. Fecha
    % 8. Autor
    
    
    \thispagestyle{empty}               % Sin encabezado ni pie de página
    \newgeometry{margin=0cm}        % Márgenes nulos para la primera página
    
    
    % Encabezado
    \fancyhead[L]{\helv #3}
    \fancyhead[R]{\helv \nouppercase{\leftmark}}
    
    
    \pagecolor{#4}        % Color de fondo para la portada
    
    \begin{figure}[p]
        \centering
        \transparent{0.3}           % Opacidad del 30% para la imagen
        
        \includegraphics[width=\paperwidth, keepaspectratio]{../../_assets/#1}
    
        \begin{tikzpicture}[remember picture, overlay]
            \node[anchor=north west, text=white, opacity=1, font=\fontsize{60}{90}\selectfont\bfseries\sffamily, align=left] at (#5, #6) {#2};
            
            \node[anchor=south east, text=white, opacity=1, font=\fontsize{12}{18}\selectfont\sffamily, align=right] at (9.7, 3) {\href{https://losdeldgiim.github.io/}{\textbf{Los Del DGIIM}, \texttt{\footnotesize losdeldgiim.github.io}}};
            
            \node[anchor=south east, text=white, opacity=1, font=\fontsize{12}{15}\selectfont\sffamily, align=right] at (9.7, 1.8) {Doble Grado en Ingeniería Informática y Matemáticas\\Universidad de Granada};
        \end{tikzpicture}
    \end{figure}
    
    
    \restoregeometry        % Restaurar márgenes normales para las páginas subsiguientes
    \pagecolor{white}       % Restaurar el color de página
    
    
    \newpage
    \thispagestyle{empty}               % Sin encabezado ni pie de página
    \begin{tikzpicture}[remember picture, overlay]
        \node[anchor=south west, inner sep=3cm] at (current page.south west) {
            \begin{minipage}{0.5\paperwidth}
                \href{https://creativecommons.org/licenses/by-nc-nd/4.0/}{
                    \includegraphics[height=2cm]{../../_assets/Licencia.png}
                }\vspace{1cm}\\
                Esta obra está bajo una
                \href{https://creativecommons.org/licenses/by-nc-nd/4.0/}{
                    Licencia Creative Commons Atribución-NoComercial-SinDerivadas 4.0 Internacional (CC BY-NC-ND 4.0).
                }\\
    
                Eres libre de compartir y redistribuir el contenido de esta obra en cualquier medio o formato, siempre y cuando des el crédito adecuado a los autores originales y no persigas fines comerciales. 
            \end{minipage}
        };
    \end{tikzpicture}
    
    
    
    % 1. Foto de fondo
    % 2. Título
    % 3. Encabezado Izquierdo
    % 4. Color de fondo
    % 5. Coord x del titulo
    % 6. Coord y del titulo
    % 7. Fecha
    % 8. Autor


}


\newcommand{\portadaBook}[8]{

    % 1. Foto de fondo
    % 2. Título
    % 3. Encabezado Izquierdo
    % 4. Color de fondo
    % 5. Coord x del titulo
    % 6. Coord y del titulo
    % 7. Fecha
    % 8. Autor

    % Personaliza el formato del título
    \pretitle{\begin{center}\bfseries\fontsize{42}{56}\selectfont}
    \posttitle{\par\end{center}\vspace{2em}}
    
    % Personaliza el formato del autor
    \preauthor{\begin{center}\Large}
    \postauthor{\par\end{center}\vfill}
    
    % Personaliza el formato de la fecha
    \predate{\begin{center}\huge}
    \postdate{\par\end{center}\vspace{2em}}
    
    \title{#2}
    \author{\href{https://losdeldgiim.github.io/}{Los Del DGIIM, \texttt{\large losdeldgiim.github.io}}
    \\ \vspace{0.5cm}#8}
    \date{Granada, #7}
    \maketitle
    
    \tableofcontents
}




\newcommand{\portadaArticle}[8]{

    % 1. Foto de fondo
    % 2. Título
    % 3. Encabezado Izquierdo
    % 4. Color de fondo
    % 5. Coord x del titulo
    % 6. Coord y del titulo
    % 7. Fecha
    % 8. Autor

    % Personaliza el formato del título
    \pretitle{\begin{center}\bfseries\fontsize{42}{56}\selectfont}
    \posttitle{\par\end{center}\vspace{2em}}
    
    % Personaliza el formato del autor
    \preauthor{\begin{center}\Large}
    \postauthor{\par\end{center}\vspace{3em}}
    
    % Personaliza el formato de la fecha
    \predate{\begin{center}\huge}
    \postdate{\par\end{center}\vspace{5em}}
    
    \title{#2}
    \author{\href{https://losdeldgiim.github.io/}{Los Del DGIIM, \texttt{\large losdeldgiim.github.io}}
    \\ \vspace{0.5cm}#8}
    \date{Granada, #7}
    \thispagestyle{empty}               % Sin encabezado ni pie de página
    \maketitle
    \vfill
}
    \portadaExamen{ffccA4.jpg}{Análisis Funcional\\Examen XVIII}{Análisis Funcional. Examen XVIII}{MidnightBlue}{-8}{28}{2026}{}

    \begin{description}
        \item[Asignatura] Análisis Funcional.
        \item[Curso Académico] 2026/2027.
        \item[Grado] Doble Grado en Ingeniería Informática y Matemáticas.
        \item[Grupo] Único.
        \item[Profesor] Miguel Martín Suárez.
        \item[Descripción] Parcial 1.
        \item[Fecha] 4 de noviembre de 2026.
        \item[Duración] 50 minutos.
    
    \end{description}
    \newpage


    % ------------------------------------

La parte práctica del Examen puntúa 2.5 puntos. Cada alumno tenía un ejercicio distinto, pero se incluyen todos para que el estudiantado tenga una mayor variedad de ejercicios para practicar.
\begin{ejercicio}[2.5 puntos]
    Consideramos el siguiente subconjunto de $C[0,1]$:
    \begin{equation*}
        X=\left\{f \in C[0,1] : f(0) + f(1) = 0\right\}
    \end{equation*}
    Se define el siguiente funcional lineal:
    \Func{\varphi}{X}{\bb{K}}{f}{\int_{0}^{1} f(t)\,dt}

    Probar que $\varphi$ es continuo, y calcular su norma. ¿Se alcanza la norma?
\end{ejercicio}

\begin{ejercicio}[2.5 puntos]
    En $X=C[0,1]$ define el siguiente funcional lineal:
    \Func{\varphi}{X}{\bb{K}}{f}{\int_{0}^{1} t^5f(t^{13})\,dt}

    Probar que $\varphi$ es continuo, y calcular su norma. ¿Se alcanza la norma?
\end{ejercicio}

\begin{ejercicio}[2.5 puntos]
    En $X=C[0,1]$ define el siguiente funcional lineal:
    \Func{\varphi}{X}{\bb{K}}{f}{\int_{0}^{\nicefrac{1}{5}} f(t)\,dt - \int_{\nicefrac{1}{3}}^{1} f(t)\,dt}

    Probar que $\varphi$ es continuo, y calcular su norma. ¿Se alcanza la norma?
\end{ejercicio}

\begin{ejercicio}[2.5 puntos]
    En $X=C[0,1]$ define el siguiente funcional lineal:
    \Func{\varphi}{X}{\bb{K}}{f}{\int_{0}^{1} (2t-1)f(t)\,dt}

    Probar que $\varphi$ es continuo, y calcular su norma. ¿Se alcanza la norma?
\end{ejercicio}

\vspace{1cm}

La parte Tipo Test del Examen puntúa 7.5 puntos. Puesto que hay 5 preguntas, cada una de ellas puntúa 1.5 puntos. Cada una de ellas tiene 4 respuestas posibles, de las cuales una, dos, tres o cuatro pueden ser correctas. La puntuación de cada pregunta se calcula de la siguiente manera:
\begin{itemize}
    \item Si se selecciona alguna respuesta incorrecta, la puntuación de la pregunta es 0.
    \item Si se seleccionan $k$ respuestas correctas y no se selecciona ninguna incorrecta, la puntuación de la pregunta es:
    \begin{equation*}
        \frac{k}{n} \cdot 1.5
    \end{equation*}
    donde $n$ es el número total de respuestas correctas de la pregunta.
\end{itemize}


Aunque cada alumno tenía un examen distinto que incluía 5 preguntas tipo test, se incluyen todas las preguntas para que el estudiantado tenga una mayor variedad de ejercicios para practicar.
\begin{ejercicio}~
    \begin{enumerate}[label=(\alph*)]
        \item Todo espacio de Banach con base de Schauder es separable.
        \item El espacio de las sucesiones acotadas, $l_\infty$, es un espacio métrico separable.
        \item $\{e_n\}_{n\in \bb{N}}$ es una base algebraica en $l_p$ para todo $1\leq p \leq \infty$ y $c_0$.
        \item Un espacio normado $X$ es separable si y solo si tiene un conjunto numerable de vectores linealmente independientes cuya expansión lineal es densa en $X$.
    \end{enumerate}    
\end{ejercicio}

\begin{ejercicio}
    Sea $Y = \bb{K}^{(\bb{N})}$ (el espacio vectorial de las sucesiones casi nulas):
    \begin{enumerate}[label=(\alph*)]
        \item La norma de la suma y la norma del máximo son equivalentes en $Y$.
        \item La norma de la suma y la norma del máximo no son equivalentes en $Y$.
        \item Cualesquiera dos normas definidas en $Y$ son equivalentes.
        \item No se puede definir ninguna norma en $Y$.
    \end{enumerate}
\end{ejercicio}

\begin{ejercicio}~
    \begin{enumerate}[label=(\alph*)]
        \item Un espacio normado es de Banach si y solo si toda serie absolutamente convergente es convergente.
        \item En un espacio normado, si $\sum\limits_{n=1}^{\infty} x_n$ es absolutamente convergente, entonces $\{x_n\}\to 0$.
        \item Si $T: X \to Y$ es un operador lineal entre espacios de Banach y $\sum\limits_{n=1}^{\infty} x_n$ es convergente, entonces $\sum\limits_{n=1}^{\infty} T(x_n)$ es convergente.
        \item Sustituir la norma de un espacio normado por otra norma equivalente puede cambiar las sucesiones de Cauchy.
    \end{enumerate}    
\end{ejercicio}

\begin{ejercicio}~
    \begin{enumerate}[label=(\alph*)]
        \item Todo operador lineal que llega a un espacio normado de dimensión finita es continuo.
        \item Si un operador lineal $T: X \to Y$ entre espacios normados no es continuo, entonces $\ker T$ es un subespacio denso de $X$.
        \item Un operador lineal $T: X \to Y$ es continuo si y solo si $\ker T$ es cerrado.
        \item Dado un espacio normado $X$ y $f\in X^*$ no nulo, existe $u\in X$ con $\|u\|=1$ tal que $|f(u)|=\|f\|$.
    \end{enumerate}    
\end{ejercicio}

\begin{ejercicio}~
    \begin{enumerate}[label=(\alph*)]
        \item Dado $1\leq p < q < \infty$, si $x\in l_p$ entonces $l_p\subset l_q$ pero $\|\cdot \|_q$ no es equivalente a $\|\cdot\|_p$.
        \item El conjunto $c_0$ de las sucesiones que tienden a cero es un subespacio cerrado de $l_\infty$.
        \item Para cualquier $1\leq p < \infty$ y todo $x\in l_p$, se tiene que $\|x\|_\infty \leq \|x\|_p$.
        \item El subespacio de las sucesiones de soporte finito $\bb{K}^{(\bb{N})}$ es denso en $l_\infty$.
    \end{enumerate}
    
\end{ejercicio}

\begin{ejercicio}~
    \begin{enumerate}[label=(\alph*)]
        \item Sea $\varphi$ una seminorma en un espacio vectorial $X$. Entonces, $\varphi(0)=0$ y $\varphi(x)\geq 0$ para todo $x\in X$.
        \item La distancia inducida por la norma $d(x,y)=\|x-y\|$ es invariante por traslaciones y homogénea por homotecias.
        \item En dimensión infinita, cualquier norma es continua pero no necesariamente uniformemente continua.
        \item Dos normas son equivalentes si y solo si existen constantes $c_1, c_2 > 0$ tales que $c_1\|x\|_1 \leq \|x\|_2 \leq c_2\|x\|_1$ para todo $x$.
    \end{enumerate}    
\end{ejercicio}




\newpage
\setcounter{ejercicio}{0}

    
\begin{ejercicio}[2.5 puntos]\label{ej:1}
    Consideramos el siguiente subconjunto de $C[0,1]$:
    \begin{equation*}
        X=\left\{f \in C[0,1] : f(0) + f(1) = 0\right\}
    \end{equation*}
    Se define el siguiente funcional lineal:
    \Func{\varphi}{X}{\bb{K}}{f}{\int_{0}^{1} f(t)\,dt}

    Probar que $\varphi$ es continuo, y calcular su norma. ¿Se alcanza la norma?\\

    Veamos que $\varphi$ es continuo:
    \begin{align*}
        |\varphi(f)| = \left|\int_{0}^{1} f(t)\,dt\right| \leq \int_{0}^{1} |f(t)|\,dt \leq \int_{0}^{1} \|f\|_\infty\,dt = \|f\|_\infty
    \end{align*}

    Por tanto, $\varphi$ es continuo y $\|\varphi\| \leq 1$. Veamos ahora que $\|\varphi\| = 1$. Para ello, para cada $n\in \bb{N}$, definimos los subconjuntos:
    \begin{align*}
        I_n^1 &= \left[0, \frac{1}{n}\right]\qquad\qquad
        I_n^2 = \left[1-\frac{1}{n}, 1\right]
    \end{align*}

    Definimos ahora la función $f_n \in C[0,1]$, mostradas en la Figura~\ref{fig:funciones_fn}, de la siguiente manera:
    \begin{equation*}
        f_n(t) = \begin{cases}
            0 & \text{si } t \in \{0, 1\}\\
            1 & \text{si } t \in [\nicefrac{1}{n}, 1 - \nicefrac{1}{n}]\\
            \text{lineal} & \text{si } t \in I_n^1 \cup I_n^2
        \end{cases}
    \end{equation*}
    \begin{figure}[h]
        \centering
        \begin{tikzpicture}
            \begin{axis}[
                axis lines = middle,
                xlabel = $t$,
                ylabel = {$f_n(t)$},
                ymin = -0.3, ymax = 1.3,
                xmin = 0, xmax = 1,
                xtick = {0, 0.5, 1},
                ytick = {-1, 0, 1},
                grid = both,
                width = 11cm,
                height = 6.5cm,
                legend pos = outer north east
            ]

            \foreach \n/\col in {3/blue, 5/red, 10/teal, 20/orange, 50/purple} {
                \edef\temp{
                    \noexpand\addplot[color=\col, thick] coordinates {
                        (0, 0)
                        (1/\n, 1)
                        (1 - 1/\n, 1)
                        (1, 0)
                    };
                    \noexpand\addlegendentry{$n = \n$}
                }
                \temp
            }

            \end{axis}
        \end{tikzpicture}
        \caption{Funciones $f_n$ para distintos valores de $n$.}
        \label{fig:funciones_fn}
    \end{figure}

    Tenemos que $f_n \in X$ y $\|f_n\|_\infty = 1$. Además, se tiene que:
    \begin{align*}
        |\varphi(f_n)| &= \left|\int_{0}^{1} f_n(t)\,dt\right| = \int_{0}^{1} f_n(t)\,dt = \int_{I_n^1} f_n(t)\,dt + \int_{\nicefrac{1}{n}}^{1 - \nicefrac{1}{n}} 1\,dt + \int_{I_n^2} f_n(t)\,dt
        \leq \\ &\leq 0+ \left(1 - \frac{2}{n}\right) + 0 = 1 - \frac{2}{n}
    \end{align*}

    Por tanto, se tiene que:
    \begin{equation*}
        \|\varphi\| = \sup\{|\varphi(f)| : f \in X, \|f\|_\infty = 1\} \geq |\varphi(f_n)| = 1 - \frac{2}{n} \quad \forall n \in \bb{N}
    \end{equation*}

    Tomando límite cuando $n \to \infty$, se tiene que:
    \begin{equation*}
        \|\varphi\| \geq 1
    \end{equation*}

    Por tanto, uniendo ambas desigualdades, se tiene que $\|\varphi\| = 1$. Veamos ahora que la norma no se alcanza. Supongamos que existe $f \in X$ con $\|f\|_\infty = 1$ tal que $|\varphi(f)| = 1$. Entonces, se tendría que:
    \begin{align*}
        1 = |\varphi(f)| = \left|\int_{0}^{1} f(t)\,dt\right| \leq \int_{0}^{1} |f(t)|\,dt \leq \int_{0}^{1} \|f\|_\infty\,dt = 1
    \end{align*}

    Por tanto, todas las desigualdades son igualdades. En particular, se tiene que:
    \begin{equation*}
        \int_{0}^{1} \|f\|_\infty\, dt = \int_{0}^{1} |f(t)|\,dt \Longrightarrow \int_{0}^{1} \left(\|f\|_\infty - |f(t)|\right)\,dt = 0
    \end{equation*}

    Como $f$ es continua y $\|f\|_\infty - |f(t)| \geq 0$ para todo $t \in [0, 1]$, se tiene para todo $t \in [0, 1]$ que:
    \begin{equation*}
        \|f\|_\infty - |f(t)| = 0 \Longrightarrow |f(t)| = \|f\|_\infty = 1
    \end{equation*}

    Por otro lado, como $\varphi(t)\in \bb{K}$, se tiene que $\exists \theta \in \bb{R}$ tal que:
    \begin{equation*}
        \varphi(f) = \int_{0}^{1} f(t)\,dt = e^{i\theta} \left|\int_{0}^{1} f(t)\,dt\right| = e^{i\theta} \int_{0}^{1} |f(t)|\,dt
    \end{equation*}

    Multiplicando por $e^{-i\theta}$, se tiene que:
    \begin{equation*}
        \int_{0}^{1} e^{-i\theta} f(t)\,dt = \int_{0}^{1} |f(t)|\,dt
    \end{equation*}

    Tomando la parte real, se tiene que:
    \begin{equation*}
        \int_{0}^{1} \Re\left(e^{-i\theta} f(t)\right)\,dt = \int_{0}^{1} |f(t)|\,dt
        \Longrightarrow
        \int_{0}^{1} \left(|f(t)| - \Re\left(e^{-i\theta} f(t)\right)\right)\,dt = 0
    \end{equation*}

    Como $|e^{-i\theta} f(t)| = |f(t)|$ y la parte real de un número complejo es menor o igual que su módulo, se tiene que:
    \begin{equation*}
        |f(t)| - \Re\left(e^{-i\theta} f(t)\right) = 0 \Longrightarrow f(t) = e^{i\theta} |f(t)|= e^{i\theta} \quad \forall t \in [0, 1]
    \end{equation*}
    
    Como $f \in X$, se tiene que:
    \begin{equation*}
        f(0) + f(1) = e^{i\theta} + e^{i\theta} = 2 e^{i\theta} = 0
    \end{equation*}
    Lo cual es una contradicción, pues $|e^{i\theta}| = 1 \neq 0$. Por tanto, la norma no se alcanza.
\end{ejercicio}


\begin{ejercicio}[2.5 puntos]
    En $X=C[0,1]$ define el siguiente funcional lineal:
    \Func{\varphi}{X}{\bb{K}}{f}{\int_{0}^{1} t^5f(t^{13})\,dt}

    Probar que $\varphi$ es continuo, y calcular su norma. ¿Se alcanza la norma?\\

    Veamos que $\varphi$ es continuo:
    \begin{align*}
        |\varphi(f)| = \left|\int_{0}^{1} t^5 f(t^{13})\,dt\right| \leq \int_{0}^{1} t^5 |f(t^{13})|\,dt \leq \int_{0}^{1} t^5 \|f\|_\infty\,dt = \frac{1}{6} \|f\|_\infty
    \end{align*}

    Por tanto, $\varphi$ es continuo y $\|\varphi\| \leq \nicefrac{1}{6}$. Veamos ahora que $\|\varphi\| = \nicefrac{1}{6}$. Para ello, consideramos la función $f_0 \in C[0,1]$ definida por:
    \begin{equation*}
        f_0(t) = 1 \quad \forall t \in [0, 1]
    \end{equation*}

    Tenemos que $\|f_0\|_\infty = 1$ y:
    \begin{align*}
        |\varphi(f_0)| = \left|\int_{0}^{1} t^5 f_0(t^{13})\,dt\right| = \int_{0}^{1} t^5\,dt = \frac{1}{6}
    \end{align*}

    Por tanto, se tiene que:
    \begin{equation*}
        \|\varphi\| = \sup\{|\varphi(f)| : f \in C[0,1], \|f\|_\infty = 1\} \geq |\varphi(f_0)| = \frac{1}{6}
    \end{equation*}

    Por tanto, uniendo ambas desigualdades, se tiene que $\|\varphi\| = \nicefrac{1}{6}$. Además, la norma se alcanza en $f_0$.
\end{ejercicio}


\begin{ejercicio}[2.5 puntos]
    En $X=C[0,1]$ define el siguiente funcional lineal:
    \Func{\varphi}{X}{\bb{K}}{f}{\int_{0}^{\nicefrac{1}{5}} f(t)\,dt - \int_{\nicefrac{1}{3}}^{1} f(t)\,dt}

    Probar que $\varphi$ es continuo, y calcular su norma. ¿Se alcanza la norma?\\

    Veamos que $\varphi$ es continuo:
    \begin{align*}
        |\varphi(f)| &= \left|\int_{0}^{\nicefrac{1}{5}} f(t)\,dt - \int_{\nicefrac{1}{3}}^{1} f(t)\,dt\right| \leq \left|\int_{0}^{\nicefrac{1}{5}} f(t)\,dt\right| + \left|\int_{\nicefrac{1}{3}}^{1} f(t)\,dt\right|
        \leq \int_{0}^{\nicefrac{1}{5}} |f(t)|\,dt + \int_{\nicefrac{1}{3}}^{1} |f(t)|\,dt \leq\\ &\leq \int_{0}^{\nicefrac{1}{5}} \|f\|_\infty\,dt + \int_{\nicefrac{1}{3}}^{1} \|f\|_\infty\,dt = \left(\frac{1}{5} - 0 + 1 - \frac{1}{3}\right) \|f\|_\infty = \frac{13}{15} \|f\|_\infty
    \end{align*}

    Por tanto, $\varphi$ es continuo y $\|\varphi\| \leq \nicefrac{13}{15}$. Veamos ahora que $\|\varphi\| = \nicefrac{13}{15}$. Para ello, la función que buscamos debería valer $1$ en $[0, \nicefrac{1}{5}]$ y $-1$ en $[\nicefrac{1}{3}, 1]$. Como $\nicefrac{1}{5} < \nicefrac{1}{3}$, podemos definir la función $f_0 \in C[0,1]$ de la siguiente manera:
    \begin{equation*}
        f_0(t) = \begin{cases}
            1 & \text{si } t \in [0, \nicefrac{1}{5}]\\
            \text{lineal} & \text{si } t \in [\nicefrac{1}{5}, \nicefrac{1}{3}]\\
            -1 & \text{si } t \in [\nicefrac{1}{3}, 1]
        \end{cases}
    \end{equation*}

    Tenemos que $\|f_0\|_\infty = 1$ y:
    \begin{align*}
        \varphi(f_0) &= \int_{0}^{\nicefrac{1}{5}} f_0(t)\,dt - \int_{\nicefrac{1}{3}}^{1} f_0(t)\,dt = \int_{0}^{\nicefrac{1}{5}} 1\,dt - \int_{\nicefrac{1}{3}}^{1} (-1)\,dt = \frac{1}{5} + \left(1 - \frac{1}{3}\right) = \frac{13}{15}
    \end{align*}

    Por tanto, $|\varphi(f_0)| = \nicefrac{13}{15}$, luego $\|\varphi\| = \nicefrac{13}{15}$. Además, la norma se alcanza en $f_0$.
\end{ejercicio}

\begin{ejercicio}[2.5 puntos]
    En $X=C[0,1]$ define el siguiente funcional lineal:
    \Func{\varphi}{X}{\bb{K}}{f}{\int_{0}^{1} (2t-1)f(t)\,dt}

    Probar que $\varphi$ es continuo, y calcular su norma. ¿Se alcanza la norma?\\

    Resuelto en el Examen XVII.
\end{ejercicio}


\begin{ejercicio}~
    \begin{enumerate}[label=(\alph*)]
        \item Todo espacio de Banach con base de Schauder es separable.
        
        Este resultado es cierto. De hecho, si $\{e_n\}_{n\in \bb{N}}$ es una base de Schauder en un espacio de Banach $X$, entonces el subespacio generado por $\{e_n\}_{n\in \bb{N}}$ es denso en $X$. Como este subespacio tiene dimensión numerable, se tiene que $X$ es separable.
        \item El espacio de las sucesiones acotadas, $l_\infty$, es un espacio métrico separable.
        
        Este resultado es falso. $\{e_n\}_{n\in \bb{N}}$ es una base de Schauder de $l_p$ para todo $1\leq p < \infty$, pero no lo es de $l_\infty$. De hecho, $l_\infty$ no es separable.

        % // TODO: Demostrado, pero fácil?
        \item $\{e_n\}_{n\in \bb{N}}$ es una base algebraica en $l_p$ para todo $1\leq p \leq \infty$ y $c_0$.
        
        No, $\{e_n\}_{n\in \bb{N}}$ es una base algebraica de las sucesiones de soporte finito $\bb{K}^{(\bb{N})}$, que es un subespacio propio de $l_p$ para todo $1\leq p \leq \infty$ y $c_0$. Por tanto, $\{e_n\}_{n\in \bb{N}}$ no es una base algebraica de $l_p$ ni de $c_0$. Notemos que, en el caso de las bases algebraicas, tan solo se consideran combinaciones lineales finitas de los vectores de la base, mientras que en el caso de las bases de Schauder se consideran combinaciones lineales infinitas.
        

        \item Un espacio normado $X$ es separable si y solo si tiene un conjunto numerable de vectores linealmente independientes cuya expansión lineal es densa en $X$.
        
        Este resultado es cierto, como se ha visto en clase.
    \end{enumerate}    
\end{ejercicio}

\begin{ejercicio}
    Sea $Y = \bb{K}^{(\bb{N})}$ (el espacio vectorial de las sucesiones casi nulas):
    \begin{enumerate}[label=(\alph*)]
        \item La norma de la suma y la norma del máximo son equivalentes en $Y$.
        
        Veamos que esto es falso. Consideremos, para cada $n\in \bb{N}$, la sucesión $x_n \in Y$ definida por:
        \begin{align*}
            x_n = \sum_{k=1}^{n} e_k
        \end{align*}

        Tenemos que $\|x_n\|_\infty = 1$ y $\|x_n\|_1 = n$. Por tanto, no existe una constante $C > 0$ tal que $\|x_n\|_\infty \leq C \|x_n\|_1$ para todo $n \in \bb{N}$. Por tanto, las normas no son equivalentes.
        \item La norma de la suma y la norma del máximo no son equivalentes en $Y$.
        
        En el apartado anterior hemos visto que esto es cierto.
        \item Cualesquiera dos normas definidas en $Y$ son equivalentes.
        
        Esto es falso. Si la dimensión de $\bb{K}^{(\bb{N})}$ fuese finita, entonces todas las normas serían equivalentes. Pero como $\bb{K}^{(\bb{N})}$ es de dimensión infinita, no todas las normas son equivalentes (como se ha visto).
        \item No se puede definir ninguna norma en $Y$.
        
        Esto es falso. Por ejemplo, la norma del máximo $\|\cdot\|_\infty$ es una norma en $Y$.
    \end{enumerate}
\end{ejercicio}

\begin{ejercicio}~
    \begin{enumerate}[label=(\alph*)]
        \item Un espacio normado es de Banach si y solo si toda serie absolutamente convergente es convergente.
        
        Este resultado es cierto.
        \item En un espacio normado, si $\sum\limits_{n=1}^{\infty} x_n$ es absolutamente convergente, entonces $\{x_n\}\to 0$.
        
        Si la serie es absolutamente convergente, entonces la serie de normas $\sum\limits_{n=1}^{\infty} \|x_n\|$ es convergente. Como el término general de una serie convergente tiende a cero, se tiene que $\|x_n\| \to 0$, y como la norma es continua, se tiene que $\{x_n\} \to 0$.
        \item Si $T: X \to Y$ es un operador lineal entre espacios de Banach y $\sum\limits_{n=1}^{\infty} x_n$ es convergente, entonces $\sum\limits_{n=1}^{\infty} T(x_n)$ es convergente.
        
        Esto es falso. Si $T$ no es continuo, no tiene por qué preservar la convergencia de las series.

        % // TODO: Contraejemplo
        \item Sustituir la norma de un espacio normado por otra norma equivalente puede cambiar las sucesiones de Cauchy.
        
        Esto es falso. Si $\|\cdot\|_1$ y $\|\cdot\|_2$ son normas equivalentes en un espacio normado $X$, entonces generan la misma topología, y por tanto las sucesiones de Cauchy son las mismas.
    \end{enumerate}    
\end{ejercicio}

\begin{ejercicio}~
    \begin{enumerate}[label=(\alph*)]
        \item Todo operador lineal que llega a un espacio normado de dimensión finita es continuo.
        
        Esto es falso. Si el dominio es de dimensión finita, entonces el operador es continuo. Si el codominio es de dimensión finita, entonces el operador es continuo si y solo si su núcleo es cerrado.
        \item Si un funcional lineal $f: X \to \mathbb{K}$ no es continuo, entonces $\ker f$ es un subespacio denso en $X$.
        
        Como la dimensión del codominio es finita, entonces el operador es continuo si y solo si su núcleo es cerrado. Por tanto, si $f$ no es continuo, entonces $\ker f$ no es cerrado. Por tanto:
        \begin{equation*}
            \ker f \subsetneq \overline{\ker f}\subset X
        \end{equation*}

        De aquí, se tiene que $\ker f$ es un subespacio denso en $X$.
        % // TODO: Por qué
        \item Un operador lineal $T: X \to Y$ entre espacios de Banach es continuo si y solo si $\ker T$ es cerrado.
        
        Esto no es cierto. Si la dimensión del codominio es finita, entonces el operador es continuo si y solo si su núcleo es cerrado; pero si la dimensión del codominio es infinita, entonces el operador puede ser discontinuo y tener el núcleo cerrado.

        % // TODO: Contraejemplo?
        \item Dado un espacio normado $X$ y $f\in X^*$ no nulo, existe $u\in X$ con $\|u\|=1$ tal que $|f(u)|=\|f\|$.
        
        Esto no es cierto, pues tan solo será cierto si $\ker f$ es proximinal, lo cual no es cierto en general. Como contraejemplo podemos emplear el funcional definido en el Ejercicio~\ref{ej:1}, que no alcanza su norma.
    \end{enumerate}    
\end{ejercicio}

\begin{ejercicio}~
    \begin{enumerate}[label=(\alph*)]
        \item Dado $1\leq p < q < \infty$, si $x\in l_p$ entonces $l_p\subset l_q$ pero $\|\cdot \|_q$ no es equivalente a $\|\cdot\|_p$.
        
        Esta afirmación es cierta.

        % // TODO: Demostración Fácil?
        \item El conjunto $c_0$ de las sucesiones que tienden a cero es un subespacio cerrado de $l_\infty$.
        
        Es cierta, puesto que $c_0$ es cerrado en $l_\infty$, que es un espacio de Banach.
        \item Para cualquier $1\leq p < \infty$ y todo $x\in l_p$, se tiene que $\|x\|_\infty \leq \|x\|_p$.
        
        Esta afirmación es cierta, pues:
        \begin{align*}
            \|x\|_\infty = \sup_{n\in \bb{N}} |x(n)| \leq \left(\sum_{n=1}^{\infty} |x(n)|^p\right)^{\nicefrac{1}{p}} = \|x\|_p
        \end{align*}
        \item El subespacio de las sucesiones de soporte finito $\bb{K}^{(\bb{N})}$ es denso en $l_\infty$.
        
        Esta afirmación es falsa. Se tiene que el cierre de $c_{00}$ (el subespacio de las sucesiones de soporte finito) en $l_\infty$ es $c_0$, que es un subespacio propio de $l_\infty$. Por tanto, $\bb{K}^{(\bb{N})}$ no es denso en $l_\infty$.
    \end{enumerate}
    
\end{ejercicio}

\begin{ejercicio}~
    \begin{enumerate}[label=(\alph*)]
        \item Sea $\varphi$ una seminorma en un espacio vectorial $X$. Entonces, $\varphi(0)=0$ y $\varphi(x)\geq 0$ para todo $x\in X$.
        
        Esta es cierta. Las condiciones de una seminorma $\varphi$ son:
        \begin{enumerate}
            \item $\varphi(x+y) \leq \varphi(x) + \varphi(y)$ para todo $x, y \in X$ (subaditividad).
            \item $\varphi(\alpha x) = |\alpha| \varphi(x)$ para todo $\alpha \in \bb{K}$ y $x \in X$ (homogeneidad).
        \end{enumerate}

        De la segunda condición, tomando $\alpha = 0$ y $x \in X$, se tiene que:
        \begin{equation*}
            \varphi(0) = \varphi(0 \cdot x) = |0| \varphi(x) = 0
        \end{equation*}

        Además, de la primera condición:
        \begin{align*}
            0=\varphi(0) = \varphi(-x + x) \leq \varphi(-x) + \varphi(x) = |-1| \varphi(x) + \varphi(x) = 2\varphi(x) \Longrightarrow\\ \Longrightarrow \varphi(x) \geq 0\qquad \forall x \in X
        \end{align*}



        \item La distancia inducida por la norma $d(x,y)=\|x-y\|$ es invariante por traslaciones y homogénea por homotecias.
        
        Esta es cierta. Veamos que $d$ es invariante por traslaciones. Fijado $z \in X$, se tiene que:
        \begin{align*}
            d(x+z, y+z) &= \| (x+z) - (y+z) \| = \| x - y \| = d(x,y)\qquad \forall x, y \in X
        \end{align*}

        Veamos que $d$ es homogénea por homotecias. Fijado $\alpha \in \bb{K}$, se tiene que:
        \begin{align*}
            d(\alpha x, \alpha y) &= \| \alpha x - \alpha y \| = \| \alpha (x - y) \| = |\alpha| \| x - y \| = |\alpha| d(x,y)\qquad \forall x, y \in X
        \end{align*}
        \item En dimensión infinita, cualquier norma es continua pero no necesariamente uniformemente continua.
        
        Esta es falsa. En cualquier espacio normado, la norma es de hecho lipschitziana con constante 1, y por tanto uniformemente continua. En particular, la desigualdad triangular inversa nos da que:
        \begin{align*}
            |\|x\| - \|y\|| \leq \|x - y\|\qquad \forall x, y \in X
        \end{align*}
        \item Dos normas son equivalentes si y solo si existen constantes $c_1, c_2 > 0$ tales que $c_1\|x\|_1 \leq \|x\|_2 \leq c_2\|x\|_1$ para todo $x$.
        
        Esta es cierta. Es una importante caracterización de la equivalencia de normas en espacios normados.
    \end{enumerate}
\end{ejercicio}




\end{document}